\SourceLeaf{413}%
\noindent Es bezeichnet dann $X,Y,Z$ die Richtung der Normale auf die Oberfläche und es ist ferner
\begin{equation}\tag{1}
\begin{aligned}
\frac{\partial x}{\partial\theta}&=m(\eta Z-\zeta Y),&\frac{\partial y}{\partial\theta}&=m(\zeta X-\xi Z),\\
\frac{\partial z}{\partial\theta}&=m(\xi Y-\eta X).
\end{aligned}
\end{equation}
\indent Man sieht hieraus, dass wenn man Einer Linie constanter $\theta$ eine bestimmte Lage im Raume angewiesen hat, die ganze Fläche eine bestimmte Lage bekommt.

Die Differentiation von (1) gibt:
\begin{equation}\tag{2}
\left\{\begin{aligned}
\frac{\partial\xi}{\partial\theta}&=\frac{\partial m}{\partial r}(\eta Z-\zeta Y)+m\left(\eta\frac{\partial Z}{\partial r}-\zeta\frac{\partial Y}{\partial r}\right)\\
\frac{\partial\eta}{\partial\theta}&=\frac{\partial m}{\partial r}(\zeta X-\xi Z)+m\left(\zeta\frac{\partial X}{\partial r}-\xi\frac{\partial Z}{\partial r}\right)\\
\frac{\partial\zeta}{\partial\theta}&=\frac{\partial m}{\partial r}(\xi Y-\eta X)+m\left(\xi\frac{\partial Y}{\partial r}-\eta\frac{\partial X}{\partial r}\right)\\
\xi.\frac{\partial m}{\partial r}&=\frac{Y\partial\zeta-Z\partial\eta}{\partial\theta}\\
\eta.\frac{\partial m}{\partial r}&=\frac{Z\partial\xi-X\partial\zeta}{\partial\theta}\\
\zeta.\frac{\partial m}{\partial r}&=\frac{X\partial\eta-Y\partial\xi}{\partial\theta}\\
\frac{\partial m}{\partial r}&=\BareRoot{\Biggl\{}\frac{\partial\xi^2+\partial\eta^2+\partial\zeta^2}{\partial\theta^2}\\
&\quad-\left(X\frac{\partial\xi}{\partial\theta}+Y\frac{\partial\eta}{\partial\theta}+Z\frac{\partial\zeta}{\partial\theta}\right)^2\Biggr\}.
\end{aligned}\right.
\end{equation}
\indent Die abermalige Differentiation gibt, $\dfrac{\partial X}{\partial r}=X'$, $\dfrac{\partial X'}{\partial r}=X''$ etc. gesetzt:
\begin{equation}\tag{3}
\left\{\begin{aligned}
X.\frac{\partial p}{\partial\theta}+p.\frac{\partial X}{\partial\theta}&=\frac{\partial\partial m}{\partial r^2}(\eta Z-\zeta Y)\\
&\quad+2\frac{\partial m}{\partial r}(\eta Z'-\zeta Y')\\
&\quad+m(\eta Z''-\zeta Y'')+mp(YZ'-ZY')\\
Y.\frac{\partial p}{\partial\theta}+p.\frac{\partial Y}{\partial\theta}&=\frac{\partial\partial m}{\partial r^2}(\zeta X-\xi Z)\\
&\quad+2\frac{\partial m}{\partial r}(\zeta X'-\xi Z')\\
&\quad+m(\zeta X''-\xi Z'')+mp(ZX'-XZ')\\
Z.\frac{\partial p}{\partial\theta}+p.\frac{\partial Z}{\partial\theta}&=\frac{\partial\partial m}{\partial r^2}(\xi Y-\eta X)\\
&\quad+2\frac{\partial m}{\partial r}(\xi Y'-\eta X')\\
&\quad+m(\xi Y''-\eta X'')+mp(XY'-YX').
\end{aligned}\right.
\end{equation}
\indent Hieraus:
\[
\begin{aligned}
\xi.\frac{\partial X}{\partial\theta}+\eta.\frac{\partial Y}{\partial\theta}+\zeta.\frac{\partial Z}{\partial\theta}
&=m(\xi YZ'+\eta ZX'+\zeta XY'\\
&\quad-\xi ZY'-\eta XZ'-\zeta YX')
\end{aligned}
\]
\noindent [oder, da]
\[
\eta Z'-\zeta Y'=qX,\quad\zeta X'-\xi Z'=qY,\quad\xi Y'-\eta X'=qZ
\]
